% Блок-схемы parse_expression() и parse_term() — рекурсивный спуск
% ГОСТ 19.701-90, src/parser.cpp

\documentclass[a4paper]{article}
\usepackage[T2A]{fontenc}
\usepackage[utf8]{inputenc}
\usepackage[russian]{babel}
\usepackage[a4paper,margin=1.5cm]{geometry}
\usepackage{tikz}
\usetikzlibrary{shapes.geometric, arrows.meta, positioning, calc}

\tikzset{
  terminal/.style={
    rectangle, rounded corners=8pt,
    draw=black, line width=1pt,
    minimum width=4.2cm, minimum height=0.85cm,
    text centered, font=\footnotesize
  },
  process/.style={
    rectangle,
    draw=black, line width=1pt,
    minimum width=4.2cm, minimum height=0.85cm,
    text centered, font=\footnotesize
  },
  decision/.style={
    diamond, aspect=2.4,
    draw=black, line width=1pt,
    inner sep=1pt,
    minimum width=3cm, minimum height=0.85cm,
    text centered, font=\footnotesize
  },
  predefined/.style={
    rectangle, double, double distance=2pt,
    draw=black, line width=1pt,
    minimum width=4.2cm, minimum height=0.85cm,
    text centered, font=\footnotesize
  },
  connector/.style={
    circle,
    draw=black, line width=1pt,
    minimum size=0.65cm,
    text centered, font=\footnotesize\bfseries
  },
  arrow/.style={->, >=Stealth, line width=1pt}
}

% ================================================================
% ЧАСТЬ 1 — parse_expression()
% Возврат цикла: одиночная стрелка слева (Z-образный путь).
% Два исхода (left+right / left-right) объединены в один блок,
% чтобы не было пересекающихся ветвей.
% ================================================================
\begin{document}
\begin{center}
\textbf{Блок-схема функции \texttt{parse\_expression()}}
\end{center}

\begin{center}
\begin{tikzpicture}[node distance=0.7cm]

\node (start) [terminal]                    {Начало parse\_expression};
\node (init)  [predefined, below=of start]  {left = parse\_term()};

\node (peek)  [predefined, below=of init]   {c = peek()};
\node (cond)  [decision,   below=of peek]   {c == '+' или \textquoteleft-\textquoteright?};

%% Выход из цикла вправо
\node (ret)   [terminal, right=4cm of cond] {Возврат left};

%% Тело цикла — вниз по главной оси
\node (eat)   [process,    below=of cond]   {eat()};
\node (peek2) [predefined, below=of eat]    {nxt = peek()};
\node (dbl)   [decision,   below=of peek2]  {nxt $\in$ \{'+', '*', '/'\}?};

%% Ошибка — вправо
\node (errDbl)[terminal, right=4cm of dbl]  {throw \textquotedbl Двойной оператор\textquotedbl};

\node (rterm) [predefined, below=of dbl]    {right = parse\_term()};
%% Один блок обновления: знак определяется сохранённым c
\node (upd)   [process,    below=of rterm]  {left = left $\pm$ right \ {\scriptsize(знак = c)}};

%% Возврат к вершине цикла: Z-путь по левому полю → прямо в peek
\coordinate (Lbot) at ($(upd.west)   + (-3.2, 0)$);
\coordinate (Ltop) at ($(peek.west)  + (-3.2, 0)$);
\draw [arrow] (upd.west) -- (Lbot) -- (Ltop) -- (peek.west);

\draw [arrow] (start)  -- (init);
\draw [arrow] (init)   -- (peek);
\draw [arrow] (peek)   -- (cond);
\draw [arrow] (cond)   -- node[above]{\footnotesize Нет} (ret);
\draw [arrow] (cond)   -- node[left] {\footnotesize Да}  (eat);
\draw [arrow] (eat)    -- (peek2);
\draw [arrow] (peek2)  -- (dbl);
\draw [arrow] (dbl)    -- node[above]{\footnotesize Да}  (errDbl);
\draw [arrow] (dbl)    -- node[left] {\footnotesize Нет} (rterm);
\draw [arrow] (rterm)  -- (upd);

\end{tikzpicture}
\end{center}

\newpage

% ================================================================
% ЧАСТЬ 2 — parse_term()
% isMul: Да -> вниз (mulL, главная ось)
%        Нет -> вправо (chkZ, errZ, divL)
% Два возврата: mulL и divL -> каждый своим Z-путём слева.
% Для того чтобы Z-пути не пересекались, они используют
% разные координаты по X: -3.2 и -4.2.
% ================================================================
\begin{center}
\textbf{Блок-схема функции \texttt{parse\_term()}}
\end{center}

\begin{center}
\begin{tikzpicture}[node distance=0.7cm]

\node (start) [terminal]                    {Начало parse\_term};
\node (init)  [predefined, below=of start]  {left = parse\_factor()};

\node (peek)  [predefined, below=of init]   {c = peek()};
\node (cond)  [decision,   below=of peek]   {c == '*' или \textquoteleft/\textquoteright?};

%% Выход из цикла вправо
\node (ret)   [terminal, right=4cm of cond] {Возврат left};

\node (eat)   [process,    below=of cond]   {eat()};
\node (peek2) [predefined, below=of eat]    {nxt = peek()};
\node (dbl)   [decision,   below=of peek2]  {nxt $\in$ \{'+', '*', '/'\}?};

\node (errDbl)[terminal, right=4cm of dbl]  {throw \textquotedbl Двойной оператор\textquotedbl};

\node (rfact) [predefined, below=of dbl]    {right = parse\_factor()};
\node (isMul) [decision,   below=of rfact]  {c == '*'?};

%% Ветка '*' — вниз по главной оси
\node (mulL)  [process,    below=1cm of isMul]  {left = left $\times$ right};

%% Ветка '/' — вправо; ниже неё проверка нуля, ещё ниже деление
\node (chkZ)  [decision,   right=4.5cm of isMul]  {right == 0?};
\node (errZ)  [terminal,   right=3cm of chkZ]      {throw \textquotedbl Деление на ноль\textquotedbl};
\node (divL)  [process,    below=1.5cm of chkZ]    {left = left $/$ right};

%% ---- стрелки ----
\draw [arrow] (start)  -- (init);
\draw [arrow] (init)   -- (peek);
\draw [arrow] (peek)   -- (cond);
\draw [arrow] (cond)   -- node[above]{\footnotesize Нет} (ret);
\draw [arrow] (cond)   -- node[left] {\footnotesize Да}  (eat);
\draw [arrow] (eat)    -- (peek2);
\draw [arrow] (peek2)  -- (dbl);
\draw [arrow] (dbl)    -- node[above]{\footnotesize Да}  (errDbl);
\draw [arrow] (dbl)    -- node[left] {\footnotesize Нет} (rfact);
\draw [arrow] (rfact)  -- (isMul);
\draw [arrow] (isMul)  -- node[left] {\footnotesize Да}  (mulL);
\draw [arrow] (isMul)  -- node[above]{\footnotesize Нет} (chkZ);
\draw [arrow] (chkZ)   -- node[above]{\footnotesize Да}  (errZ);
\draw [arrow] (chkZ)   -- node[left] {\footnotesize Нет} (divL);

%% Возврат mulL → peek: Z-путь по левому полю, вход слева
\coordinate (LbotM) at ($(mulL.west)  + (-3.2, 0)$);
\coordinate (LtopM) at ($(LbotM      |- peek.west)$);
\draw [arrow] (mulL.west) -- (LbotM) -- (LtopM) -- (peek.west);

%% Возврат divL → peek: Z-путь по левому полю, вход сверху
\coordinate (LbotD) at ($(divL.west)  + (-9.0, 0)$);
\coordinate (LtopD) at ($(LbotD      |- peek.north)$);
\draw [arrow] (divL.west) -- (LbotD) -- (LtopD) -- (peek.north);

\end{tikzpicture}
\end{center}

\end{document}
